% ============================================================
% SECCIÓN 1: Introducción y Visión del Proyecto
% ============================================================
\chapter*{Introducción y Visión del Proyecto}
\addcontentsline{toc}{chapter}{Introducción y Visión del Proyecto}

% ─────────────────────────────────────────────────────────────
\section{Origen del Proyecto}

El \textbf{Proyecto Demeter} nace en el seno de la asociación universitaria
\textbf{ESIBot} (Escuela Técnica Superior de Ingeniería, Universidad de Sevilla)
como un ejercicio de ingeniería full-stack aplicada a un dominio físico real:
la monitorización y el control automatizado de entornos de cultivo. Su autor,
Daniel Montero Fernández, conduce el proyecto en solitario a lo largo de
varias iteraciones, partiendo de una primera prueba de concepto sobre Arduino
y evolucionando hasta la arquitectura distribuida documentada en este texto:
una red de microcontroladores ESP32, un servicio de borde sobre Raspberry Pi,
un servidor central desplegable indistintamente en Google Cloud Platform o en
hardware doméstico, y un SDK científico de Python para análisis agronómico.

La presente versión 1.0 marca el cierre del ciclo de desarrollo activo del
proyecto. Se entrega como \textit{release} estable y \textit{portfolio
piece}, no como abandono: el sistema queda funcional, documentado y
reproducible, con una cláusula explícita de no continuidad (\textit{no further
development planned}) que libera al autor del compromiso de mantenimiento
indefinido y, simultáneamente, fija el alcance de lo entregado.

% ─────────────────────────────────────────────────────────────
\section{Motivación}

La monitorización continua de variables ambientales y la trazabilidad de los
sujetos experimentales son requisitos básicos en cualquier protocolo
agronómico riguroso, pero su implantación en pequeños invernaderos de
investigación universitaria choca con dos limitaciones recurrentes: la
\textbf{necesidad de presencia humana} para tomar lecturas y operar el riego,
y el \textbf{coste de las plataformas comerciales} de gestión de información
de laboratorio (LIMS), normalmente diseñadas para entornos clínicos o
industriales y desproporcionadas para grupos pequeños. Demeter aborda ambos
frentes con hardware de bajo coste, software libre y un alcance
deliberadamente acotado.

% ─────────────────────────────────────────────────────────────
\section{Objetivos del Sistema}

\subsection{Objetivo General}
Desarrollar una plataforma IoT integral que permita monitorizar, controlar y
analizar entornos de cultivo, proporcionando trazabilidad científica completa
de las plantas y los experimentos asociados.

\subsection{Objetivos Específicos}

Estos cuatro objetivos articulan el resto del documento y se retoman en el
Capítulo~\ref{sec:conclusiones} para evaluar el grado de cumplimiento de
la versión 1.0:

\begin{itemize}[leftmargin=*, label=--]
    \item \textbf{Automatización IoT.} Desplegar una red de sensores de bajo
          coste y alta fiabilidad que recolecte variables ambientales de
          forma continua y sin intervención humana.
    \item \textbf{Control Remoto.} Permitir la actuación sobre los elementos
          físicos del entorno (bombas de riego, electroválvulas, ventanas)
          desde cualquier punto con acceso a internet.
    \item \textbf{Trazabilidad LIMS.} Implementar un sistema de gestión de
          información de laboratorio que mantenga un registro estructurado
          de cada planta y los experimentos en los que participa.
    \item \textbf{Análisis Científico.} Dotar a los investigadores de
          herramientas (SDK, gráficas, informes exportables) para extraer y
          analizar datos sin depender del equipo de soporte técnico.
\end{itemize}

% ─────────────────────────────────────────────────────────────
\section{Alcance del Documento}

Este documento es la \textbf{documentación técnica de arquitectura} del
sistema Demeter. Su audiencia objetivo son personas con perfil técnico
(desarrolladores, evaluadores, mantenedores futuros) que necesitan entender
\textit{cómo está construido} el sistema, \textit{por qué se tomaron las
decisiones} de diseño que se tomaron, y \textit{cómo desplegarlo} desde un
entorno limpio.

\textbf{Lo que cubre este documento:}

\begin{itemize}[leftmargin=*, label=--]
    \item Arquitectura de cada una de las cuatro capas del sistema y de las
          fronteras entre ellas (Capítulos~\ref{sec:protocolo} a \ref{sec:sdk}).
    \item Estrategia de validación y banco de tests automatizados verificable
          (Capítulo~\ref{sec:testing}).
    \item Procedimiento reproducible de despliegue en los cuatro entornos
          soportados, incluyendo CI/CD y operaciones
          (Capítulo~\ref{sec:despliegue}).
    \item Estado de la versión 1.0 y lecciones aprendidas
          (Capítulo~\ref{sec:conclusiones}).
\end{itemize}

\textbf{Lo que NO cubre este documento} (y dónde encontrarlo):

\begin{itemize}[leftmargin=*, label=--]
    \item \textbf{Manual de usuario final}: cómo usar el dashboard, cómo crear
          un experimento, cómo interpretar las gráficas. La interfaz se
          documenta en el propio frontend mediante \textit{tooltips}, ayudas
          contextuales y un README específico del módulo.
    \item \textbf{Informe de viabilidad y cierre del proyecto}: análisis de
          sostenibilidad económica, justificación del cierre como v1.0
          \textit{release} y plan de entrega académica. Documento separado en
          \texttt{docs/Closure/LaTeX/main.pdf}.
    \item \textbf{Guías de desarrollo de plugins o módulos terceros}: la
          versión 1.0 no contempla extensibilidad mediante plugins. Una
          eventual versión futura debería añadir esta dimensión a la
          documentación.
\end{itemize}

% ─────────────────────────────────────────────────────────────
\section{Glosario y Siglas}

A lo largo del documento se emplean los siguientes términos y siglas, muchos
de los cuales son habituales en sus dominios respectivos pero pueden requerir
contextualización conjunta.

\begin{table}[H]
    \centering
    \caption{Glosario de términos y siglas empleados en el documento}
    \begin{tabularx}{\textwidth}{lX}
        \toprule
        \textbf{Sigla / Término}      & \textbf{Significado y contexto en Demeter} \\
        \midrule
        ADR                            & \textit{Architecture Decision Record}. Registro corto que documenta una decisión arquitectónica relevante, sus alternativas y su justificación. \\
        API REST                       & \textit{Application Programming Interface} sobre HTTP siguiendo el estilo arquitectónico REST. \\
        CRC                            & \textit{Cyclic Redundancy Check}. Suma de verificación que protege la integridad de las tramas del protocolo binario Demeter. \\
        ESP-NOW                        & Protocolo de comunicación inalámbrica entre dispositivos ESP32 sin necesidad de un punto de acceso WiFi. \\
        ET\textsubscript{0}            & \textit{Evapotranspiración de referencia}. Indicador agronómico que estima la pérdida de agua del cultivo por unidad de tiempo. \\
        GDD                            & \textit{Growing Degree Days}. Acumulado térmico que se usa para predecir fases fenológicas del cultivo. \\
        IoT                            & \textit{Internet of Things}. Paradigma de dispositivos físicos conectados que recolectan y reaccionan a datos del entorno. \\
        JWT                            & \textit{JSON Web Token}. Formato compacto de token firmado usado para autenticación entre el frontend y el backend. \\
        LIMS                           & \textit{Laboratory Information Management System}. Sistema que organiza la trazabilidad de muestras y experimentos en un laboratorio. \\
        OTA                            & \textit{Over-The-Air}. Actualización remota del firmware sin intervención física sobre el dispositivo. NO implementado en v1.0. \\
        Parquet                        & Formato columnar de almacenamiento de datos usado por el SDK como caché local de telemetría. \\
        RBAC                           & \textit{Role-Based Access Control}. Modelo de control de acceso basado en roles (admin, viewer) usado por el backend. \\
        RQ                             & \textit{Redis Queue}. Librería Python que implementa una cola de tareas asíncronas sobre Redis, usada por el \textit{math worker}. \\
        SDK                            & \textit{Software Development Kit}. En Demeter, el paquete Python \texttt{demeter\_sdk} que provee acceso programático a la plataforma. \\
        TLS                            & \textit{Transport Layer Security}. Cifrado del transporte HTTPS y WebSocket Seguro (WSS). \\
        UART                           & \textit{Universal Asynchronous Receiver-Transmitter}. Enlace serie entre el ESP32 \textit{gateway} y la Raspberry Pi. \\
        VPD                            & \textit{Vapor Pressure Deficit}. Indicador agronómico que cuantifica el estrés hídrico atmosférico sobre la planta. \\
        WSS                            & \textit{WebSocket Secure}. WebSocket sobre TLS, transporte usado entre la Raspberry Pi y el backend. \\
        \bottomrule
    \end{tabularx}
    \label{tab:glosario}
\end{table}
